\documentclass[aspectratio=169,10pt]{beamer}
\usepackage{lmodern}
\renewcommand{\familydefault}{\sfdefault}
\usepackage{amsmath, amsfonts}
\usepackage{graphicx}
\usepackage{booktabs}
\usepackage{enumitem}
\usepackage{tcolorbox}

% Layout
\setbeamertemplate{navigation symbols}{}
\setbeamercolor{frametitle}{fg=blue!70!black}
\setbeamercolor{title}{fg=blue!70!black}
\setbeamerfont{frametitle}{series=\bfseries}
\setlist[itemize]{noitemsep,topsep=2pt,partopsep=0pt,parsep=0pt,leftmargin=*}
\setlength{\abovedisplayskip}{4pt}
\setlength{\belowdisplayskip}{4pt}

% Header/footer border lines
\setbeamertemplate{headline}{\rule{\paperwidth}{0.3pt}}
\setbeamertemplate{footline}{\rule{\paperwidth}{0.3pt}}

% Box for key equations
\newtcolorbox{keyeq}{
  colback=blue!3,
  colframe=blue!50!black,
  boxrule=0.6pt,
  arc=2pt,
  left=6pt,
  right=6pt,
  top=4pt,
  bottom=4pt
}

\title{Linear Regression: Core Concepts and Diagnostics}
\author{Thomaskutty Reji}
\date{}

\begin{document}

% ------------------------------------------------------------
\begin{frame}[t]{Simple Linear Regression}
\textbf{What it is}
\begin{itemize}
\item Models a linear relationship between a response $y$ and a single predictor $x$.
\item Goal: estimate the line that minimizes squared errors.
\end{itemize}

\begin{keyeq}
\[
\hat{y} = \beta_0 + \beta_1 x
\]
\end{keyeq}

\textbf{Least squares objective}
\[
\min_{\beta_0,\beta_1} \sum_{i=1}^{n} (y_i - (\beta_0 + \beta_1 x_i))^2
\]
\end{frame}

% ------------------------------------------------------------
\begin{frame}[t]{Multiple Linear Regression (Matrix Form)}
\textbf{Formulation}
\begin{keyeq}
\[
\hat{\beta} = (X^\top X)^{-1} X^\top y
\]
\end{keyeq}
\small
$X \in \mathbb{R}^{n\times p}$, $y \in \mathbb{R}^{n}$, $\hat{\beta} \in \mathbb{R}^{p}$.

\textbf{Algorithm (OLS)}
\begin{itemize}
\item Construct $X$ with intercept column.
\item Compute $X^\top X$ and $X^\top y$.
\item Solve for $\hat{\beta}$ (normal equations).
\end{itemize}
\end{frame}

% ------------------------------------------------------------
\begin{frame}[t]{Regularization Techniques}
\textbf{Why regularize}
\begin{itemize}
\item Reduce variance and improve generalization.
\item Stabilize estimates under multicollinearity.
\end{itemize}

\textbf{Ridge (L2)}
\[
\hat{\beta}_{ridge}=\arg\min_{\beta}\; \|y-X\beta\|_2^2 + \lambda\|\beta\|_2^2
\]

\textbf{Lasso (L1)}
\[
\hat{\beta}_{lasso}=\arg\min_{\beta}\; \|y-X\beta\|_2^2 + \lambda\|\beta\|_1
\]

\textbf{Elastic Net}
\[
\hat{\beta}_{EN}=\arg\min_{\beta}\; \|y-X\beta\|_2^2 + \lambda\left(\alpha\|\beta\|_1 + (1-\alpha)\|\beta\|_2^2\right)
\]
\end{frame}

% ------------------------------------------------------------
\begin{frame}[t]{Error Metrics and Fit Quality}
\textbf{Residual Standard Error (RSE)}
\[
RSE = \sqrt{\frac{1}{n-p-1} \sum_{i=1}^{n} (y_i-\hat{y}_i)^2}
\]

\textbf{R\textsuperscript{2}}
\[
R^2 = 1 - \frac{RSS}{TSS}
\]

\textbf{Adjusted R\textsuperscript{2}}
\[
\text{Adj } R^2 = 1 - \frac{(1-R^2)(N-1)}{N-p-1}
\]

\textbf{AIC / BIC (model selection)}
\begin{itemize}
\item Penalize model complexity to avoid overfitting.
\end{itemize}
\end{frame}

% ------------------------------------------------------------
\begin{frame}[t]{Correlation and Covariance}
\textbf{Correlation}
\[
\text{Cor}(X,Y) = \frac{\sum_{i=1}^{n}(x_i-\bar{x})(y_i-\bar{y})}{\sqrt{\sum(x_i-\bar{x})^2 \sum(y_i-\bar{y})^2}}
\]

\textbf{Covariance}
\[
\text{Cov}(X,Y) = \frac{\sum (X_i-\bar{X})(Y_i-\bar{Y})}{N}
\]
\end{frame}

% ------------------------------------------------------------
\begin{frame}[t]{Hypothesis Testing in Regression}
\textbf{t-test for a coefficient}
\begin{itemize}
\item $H_0: \beta_1 = 0$ (no relationship)
\item $H_1: \beta_1 \neq 0$
\end{itemize}
\[
 t = \frac{\hat{\beta}_1 - 0}{SE(\hat{\beta}_1)}
\]

\textbf{F-test (overall significance)}
\begin{itemize}
\item $H_0: \beta_1 = \beta_2 = \cdots = \beta_p = 0$
\item $H_a:$ at least one $\beta_j \neq 0$
\end{itemize}
\end{frame}

% ------------------------------------------------------------
\begin{frame}[t]{Confidence and Prediction Intervals}
\textbf{Confidence interval}
\begin{itemize}
\item Range likely to contain the true mean response.
\end{itemize}

\textbf{Prediction interval}
\begin{itemize}
\item Range likely to contain a new observation.
\end{itemize}
\end{frame}

% ------------------------------------------------------------
\begin{frame}[t]{OLS Assumptions (Violations)}
\begin{itemize}
\item Linearity
\item Independence of errors
\item Homoscedasticity
\item Normality of errors
\item No (severe) multicollinearity
\end{itemize}
\end{frame}

% ------------------------------------------------------------
\begin{frame}[t]{Residual Diagnostics}
\begin{itemize}
\item Residuals vs fitted
\item Q-Q plot
\item Shapiro-Wilk, KS, Anderson-Darling
\item One-sample t-test on residual mean
\item Leverage and Cook's distance
\item Scale-Location plot
\item Breusch-Pagan (heteroscedasticity)
\item Durbin-Watson (autocorrelation)
\item VIF (multicollinearity)
\end{itemize}
\end{frame}

\end{document}
